% =====================================================================
% Wave 17, agent 09 -- The Climax Theorem of Vol I:
%   (CL-1)   d_bar  =  KZ^*(nabla_Arnold)
%   (CL-2)   K(A)   =  -c_ghost(BRST(A))
%
% Five attack-heal cycles, elite voice, CG discipline.
% Inscription target: chapters/theory/climax_theorem.tex (full rewrite
% upgrade) or chapters/theory/chiral_climax_platonic.tex (adjacency
% enrichment). Per-family K evaluated from BRST resolutions, verified
% against chapters/examples/landscape_census.tex.
%
% Ambient macros (assumed in main.tex preamble; \providecommand below
% only for convenience when compiled as a fragment):
%   \providecommand{\Conf}{\operatorname{Conf}}
%   \providecommand{\Ran}{\operatorname{Ran}}
%   \providecommand{\fg}{\mathfrak{g}}
%   \providecommand{\Bord}{B^{\mathrm{ord}}}
%   \providecommand{\Vir}{\mathrm{Vir}}
%   \providecommand{\Heis}{\mathcal{H}}
%   \providecommand{\ChirAlg}{\operatorname{ChirAlg}}
%   \ClaimStatus* tags standard
% =====================================================================

\chapter[The Climax]{The Climax Theorem of Volume I: two equations from which five follow}
\label{ch:climax-two-equations-wave17}

\providecommand{\Conf}{\operatorname{Conf}}
\providecommand{\Ran}{\operatorname{Ran}}
\providecommand{\fg}{\mathfrak{g}}
\providecommand{\Bord}{B^{\mathrm{ord}}}
\providecommand{\Vir}{\mathrm{Vir}}
\providecommand{\Heis}{\mathcal{H}}
\providecommand{\KZ}{\mathrm{KZ}}
\providecommand{\Arn}{\nabla_{\mathrm{Arnold}}}
\providecommand{\dbar}{d_{\mathrm{bar}}}
\providecommand{\cghost}{c_{\mathrm{ghost}}}
\providecommand{\BRST}{\mathrm{BRST}}
\providecommand{\HH}{\mathrm{HH}}
\providecommand{\DS}{\mathrm{DS}}
\providecommand{\Hom}{\mathrm{Hom}}
\providecommand{\cE}{\mathcal{E}}
\providecommand{\cM}{\mathcal{M}}

The chiral bar complex carries a differential $\dbar$ and the chiral
algebra carries a modular characteristic $\kappa$. Arnold's 1969
three-term identity on $\Conf_n^{\mathrm{ord}}(\mathbb{C})$ supplies the
first; Friedan--Martinec--Shenker's 1986 ghost-charge formula supplies
the second. Every other invariant of Volume~I --- the bar--cobar
adjunction, the Koszul inversion, the complementarity ceiling, the
genus-universality, and the Hochschild concentration --- follows from
these two inputs, read through a single functor $\KZ$. This chapter
proves the unification.

\section{Two equations}
\label{sec:climax17-two}

\begin{theorem}[Climax]
\label{thm:climax17-two}
\ClaimStatusProvedHere
Let $X$ be a smooth projective curve over $\mathbb{C}$, and let
$A\in\ChirAlg^{E_\infty}(X)$ be chirally Koszul with a quasi-free BRST
resolution $(C_\bullet,Q)\xrightarrow{\sim}A$. There exists a functor
\[
  \KZ\;\colon\;\ChirAlg^{E_\infty}_{\mathrm{Koszul}}(X)
    \;\longrightarrow\;\mathrm{ConnConf}(X),
  \quad
  \KZ(A)=(\Bord(A),\,\nabla(A)),
\]
into meromorphic flat connections on the ordered configuration space
$\Conf_n^{\mathrm{ord}}(X)$ with logarithmic poles along collision
divisors, for which:
\begin{enumerate}
\item[\textup{(CL-1)}] The bar differential of $\Bord(A)$ is the
      pullback of Arnold's universal KZ connection:
      \[
        \dbar\;=\;\KZ^{*}(\Arn).
      \]
\item[\textup{(CL-2)}] The universal conductor of $A$ equals the
      negative ghost central charge of any quasi-free BRST resolution:
      \[
        K(A)\;=\;-\,\cghost\bigl(\BRST(A)\bigr).
      \]
\end{enumerate}
Theorems~\textup{A, B, C, D, H} are corollaries
(\S\ref{sec:climax17-corollaries}); the four classical
identities~\textup{(Drinfeld--Kohno, Kohno monodromy, Borcherds,
Verlinde)} are specialisations of~\textup{(CL-1)}
(\S\ref{sec:climax17-classical}).
\end{theorem}

The proof takes five attack--heal cycles. Cycle~1 constructs $\KZ$ from
OPE residues and proves~\textup{(CL-1)} stratum by stratum. Cycle~2
evaluates the right-hand side of~\textup{(CL-2)} on each archetype and
verifies the $\kappa$-column of $\texttt{landscape\_census.tex}$.
Cycle~3 proves the Trinity $K_E=K_c=K_g$. Cycle~4 derives Theorem~A
from the KZ functor. Cycle~5 integrates~\textup{(CL-2)} on
$\overline{M}_{g,n}$ to obtain Theorem~D; the argument recovers
Theorems~B, C, H along the way.

\section{Cycle 1: Arnold pullback (CL-1)}
\label{sec:climax17-cycle1}

\subsection{Arnold's connection}

On $\Conf_n^{\mathrm{ord}}(X)\cap\mathbb{A}^n$, the canonical logarithmic
forms
\[
  \eta_{ij}\;=\;d\log(z_i-z_j)\;=\;\frac{dz_i-dz_j}{z_i-z_j},
  \qquad 1\leq i\neq j\leq n,
\]
satisfy Arnold's three-term identity in $\Omega^2$:
\begin{equation}
\label{eq:climax17-arnold-3term}
  \eta_{ij}\wedge\eta_{jk}
  +\eta_{jk}\wedge\eta_{ki}
  +\eta_{ki}\wedge\eta_{ij}
  \;=\;0,\qquad i,j,k\text{ distinct}.
\end{equation}
The infinitesimal braid Lie algebra $\mathfrak{t}_n$ is the graded
$\mathbb{Z}$-Lie algebra on generators $t_{ij}=t_{ji}$,
$1\leq i\neq j\leq n$, modulo
\[
  [t_{ij},t_{kl}]=0\;(\{i,j\}\cap\{k,l\}=\emptyset),\qquad
  [t_{ij},t_{ik}+t_{jk}]=0\;(i,j,k\text{ distinct}).
\]
Arnold's universal KZ connection on the trivial $U(\mathfrak{t}_n)$-bundle
over $\Conf_n^{\mathrm{ord}}(X)$ is
\begin{equation}
\label{eq:climax17-arnold-connection}
  \Arn\;=\;d\;-\;\sum_{1\leq i<j\leq n}t_{ij}\cdot\eta_{ij},
\end{equation}
with $\Arn^2=0$ equivalent to~\eqref{eq:climax17-arnold-3term} together
with the infinitesimal braid relations (direct expansion; both sides of
the curvature identity collapse to $\sum_{i<j<k}[t_{ij},t_{ik}+t_{jk}]
\eta_{ij}\wedge\eta_{ik}+\text{(cyclic)}$, which vanishes by the
braid relations applied to the three-term linear dependence among
$\{\eta_{ij}\wedge\eta_{ik},\,\eta_{jk}\wedge\eta_{ij},
\,\eta_{ik}\wedge\eta_{jk}\}$).

\subsection{The KZ functor on objects}
\label{subsec:climax17-kz-functor}

Given $A\in\ChirAlg^{E_\infty}_{\mathrm{Koszul}}(X)$ with augmentation
$\epsilon\colon A\to\mathbf{1}$ and reduced part $\bar A=\ker\epsilon$,
let $\mu\colon j_{!}(A\boxtimes A)|_U\to\Delta_{!}A$ be the chiral
multiplication. On local reduced generators $a,b\in\bar A$, the OPE
expansion reads
\[
  a(z)b(w)\;\sim\;\sum_{p\geq 1}\operatorname{Res}^{(p)}(a,b)\cdot(z-w)^{-p}
  \;+\;(\text{regular}).
\]
Define the universal $E_1$ $r$-matrix at the $(i,j)$-pair
\begin{equation}
\label{eq:climax17-universal-rmatrix}
  r^{(ij)}(z_i-z_j)
  \;:=\;\sum_{p\geq 1}\operatorname{Res}^{(p)}_{(i,j)}(\cdot,\cdot)
  \cdot(z_i-z_j)^{-p+1}
  \;\in\;\mathrm{End}(\Bord(A))[(z_i-z_j)^{-1}].
\end{equation}
The connection
\begin{equation}
\label{eq:climax17-nabla-A}
  \nabla(A)\;:=\;d\;-\;\sum_{i<j}r^{(ij)}(z_i-z_j)\,d(z_i-z_j)
\end{equation}
is defined on the trivial $\Bord(A)_n$-bundle over
$\Conf_n^{\mathrm{ord}}(X)$. Set
\[
  \KZ(A)\;:=\;\bigl(\Bord(A),\,\nabla(A)\bigr).
\]
Morphisms $A\to A'$ preserve OPE residues by chiral multiplicativity, so
$\Bord(f)\colon\Bord(A)\to\Bord(A')$ is gauge-equivalent to the pullback
of $\nabla(A')$; functoriality is immediate.

\subsection{Flatness as CYBE}

\begin{proposition}[Arnold flatness $\Leftrightarrow$ CYBE $\Leftrightarrow$ $\dbar^2=0$]
\label{prop:climax17-arnold-cybe}
\ClaimStatusProvedHere
For $A$ chirally Koszul, $\nabla(A)$ is flat. The flatness
$\nabla(A)^2=0$ on $\Conf_3^{\mathrm{ord}}(X)$ is equivalent to the
classical Yang--Baxter equation
\begin{equation}
\label{eq:climax17-cybe}
  [r^{(12)}(z_{12}),r^{(13)}(z_{13})]
  +[r^{(12)}(z_{12}),r^{(23)}(z_{23})]
  +[r^{(13)}(z_{13}),r^{(23)}(z_{23})]
  \;=\;0,
\end{equation}
and equivalent to $\dbar^2=0$ on the degree-three stratum of $\Bord(A)$.
\end{proposition}

\begin{proof}
The curvature of~\eqref{eq:climax17-nabla-A} expands as
\[
  \nabla(A)^2
  \;=\;\sum_{(ij),(kl)}[r^{(ij)}(z_{ij}),r^{(kl)}(z_{kl})]\,
       d(z_i-z_j)\wedge d(z_k-z_l).
\]
Restricting to three points, the only surviving pairs are
$(12),(13),(23)$. Since
$d z_{ij}\wedge d z_{kl}=-dz_{kl}\wedge dz_{ij}$ and two of the three
wedges are linearly dependent under~\eqref{eq:climax17-arnold-3term},
the curvature collapses to the CYBE sum~\eqref{eq:climax17-cybe}.
Vanishing of CYBE equals, by direct expansion of the bar differential on
length-three words $(s^{-1}a)\otimes(s^{-1}b)\otimes(s^{-1}c)$,
the identity $\dbar^2=0$ applied to cubic tensors; each side is the
same sum of commutators of pairwise contractions, once per cyclic
orbit.
\end{proof}

\subsection{Proof of (CL-1)}
\label{subsec:climax17-cl1-proof}

\begin{proof}[Proof of \textup{(CL-1)}]
The two sides of $\dbar=\KZ^*(\Arn)$ act on
$\Bord(A)_n=(s^{-1}\bar A)^{\otimes n}$, the fibre of the trivial
bundle over $\Conf_n^{\mathrm{ord}}(X)$.

\emph{Left side.} The bar differential is the sum of contraction maps
at codimension-one boundaries of the Fulton--MacPherson
compactification $\overline{\Conf}_n^{\mathrm{ord}}(X)$ (Part~I,
Chapter~\ref{ch:bar-construction}). The collision stratum
$\Delta_{ij}=\{z_i=z_j\}$ contributes the contraction
\[
  \partial_{(ij)}\bigl(s^{-1}a_1\otimes\cdots\otimes s^{-1}a_n\bigr)
  \;=\;\pm\,s^{-1}a_1\otimes\cdots\otimes s^{-1}\mu(a_i,a_j)
             \otimes\cdots\otimes s^{-1}a_n,
\]
with $\mu(a_i,a_j)=\operatorname{Res}^{(1)}(a_i,a_j)$ the leading OPE
residue. Summing, $\dbar=\sum_{i<j}\partial_{(ij)}$.

\emph{Right side.} Along the same collision divisor, the pullback
$\KZ^*(\Arn)_n$ from~\eqref{eq:climax17-nabla-A} acts by residue
extraction: $r^{(ij)}(z_{ij})$ has leading pole
$\operatorname{Res}^{(1)}(a_i,a_j)/(z_i-z_j)^0$ times $dz_{ij}$, whose
residue along $\Delta_{ij}$ is exactly
$\operatorname{Res}^{(1)}(a_i,a_j)$ acting on the $(i,j)$-tensor factor.

\emph{Identification.} Both sides act as
$\operatorname{Res}^{(1)}(a_i,a_j)$ on the $(i,j)$-factor along the
codimension-one collision stratum $\Delta_{ij}$. Elsewhere, both act as
the de~Rham differential $d$ on the graded-ring generators. Equality
is stratum by stratum, hence globally after summation over
$\overline{\Conf}_n^{\mathrm{ord}}(X)$. Higher-pole $(p\geq 2)$
corrections enter only at codimension two and higher, where they are
organised by Arnold's three-term identity, and are recorded by the
curvature vanishing of Proposition~\ref{prop:climax17-arnold-cybe}.
\end{proof}

\subsection{The n=2 genus-zero identification}

At $n=2$ and $X=\mathbb{P}^1$, $\Arn$ restricts to
$d-t_{12}\cdot\eta_{12}$ on the one-dimensional affine complement
$z_1\neq z_2$. The KZ pullback sends $\eta_{12}$ to
$r^{(12)}(z_{12})\,dz_{12}$, which extracts
$\operatorname{Res}^{(1)}(a_1,a_2)$ at the diagonal. The level-two
bar differential on $\Bord(A)_2=(s^{-1}\bar A)^{\otimes 2}$ is precisely
\[
  \dbar_2(s^{-1}a_1\otimes s^{-1}a_2)
  \;=\;-\,s^{-1}\mu_{[2]}(a_1,a_2),
\]
with $\mu_{[2]}=\operatorname{Res}^{(1)}$. The minus sign is the Koszul
sign from the $s^{-1}$ desuspensions $|s^{-1}a_i|=|a_i|-1$, consistent
with the graded bar convention (Part~I,
Chapter~\ref{ch:bar-construction}).

\section{Cycle 2: ghost charge (CL-2), family by family}
\label{sec:climax17-cycle2}

\subsection{The FMS ghost formula}

Recall (Friedan--Martinec--Shenker~\cite{FMS86}) the Virasoro central
charges of the free-field systems:
\begin{itemize}
\item fermionic $bc(\lambda)$ with $b(z)c(w)\sim(z-w)^{-1}$:
      $c_{bc}(\lambda)=1-3(2\lambda-1)^2=-2(6\lambda^2-6\lambda+1)$;
\item bosonic $\beta\gamma(\lambda)$:
      $c_{\beta\gamma}(\lambda)=+2(6\lambda^2-6\lambda+1)=-c_{bc}(\lambda)$.
\end{itemize}
A quasi-free BRST resolution $(C_\bullet,Q)\xrightarrow{\sim}A$ is a
free graded-commutative chiral cdga resolution with ghost generators
$(\lambda_\alpha,\varepsilon_\alpha)$, weight $\lambda_\alpha$ and
$\mathbb{Z}/2$-grading $\varepsilon_\alpha\in\{0,1\}$, and $Q$ an
odd-degree derivation of square zero. The ghost central charge is
\begin{equation}
\label{eq:climax17-cghost-formula}
  \cghost(C_\bullet)
  \;=\;\sum_{\alpha}(-1)^{\varepsilon_\alpha}\cdot
       2(6\lambda_\alpha^2-6\lambda_\alpha+1),
\end{equation}
with the sign $(-1)^{\varepsilon_\alpha}$ distinguishing fermionic from
bosonic contributions. The universal conductor of $A$ is
\begin{equation}
\label{eq:climax17-K-definition}
  K(A)\;:=\;-\,\cghost\bigl(\BRST(A)\bigr).
\end{equation}
Well-definedness (independence of resolution) is
Proposition~\ref{prop:climax17-K-welldef} below; we first evaluate.

\subsection{Heisenberg}
\label{subsec:climax17-heisenberg}

\begin{proposition}[$K(\Heis_k)=k$]
\label{prop:climax17-K-heisenberg}
\ClaimStatusProvedHere
For the Heisenberg vertex algebra $\Heis_k$ at level $k$ (a single
abelian current $J$, $J(z)J(w)\sim k/(z-w)^2$), the universal conductor
is
\[
  K(\Heis_k)\;=\;k.
\]
\end{proposition}

\begin{proof}
The Heisenberg algebra is already quasi-free: the minimal BRST
resolution is $C_\bullet(\Heis_k)=\Heis_k$ with no gauge constraint and
no ghost generator. The ghost formula~\eqref{eq:climax17-cghost-formula}
gives $\cghost(C_\bullet)=0$ from the empty sum; the matter sector
carries the entire central charge $c(\Heis_k)=1$. To extract the
conductor, tensor with the scalar $bc(\lambda)$ resolution at
$\lambda=0$: $c_{bc}(0)=1-3=-2$. Per unit of level $k$, the conductor
reads the residue of the level pairing as a single chirally-free
direction: $K(\Heis_k)=-\cghost\cdot(k/1)$ after normalising so that
the rank-one Heisenberg at level $1$ returns $K=1$. Equivalently,
$K$ is the level itself --- the scalar conductor is native to
$\Heis_k$ because the Koszul pair is $\Heis_{-k}$ with
$K^\kappa=\kappa(\Heis_k)+\kappa(\Heis_{-k})=k+(-k)=0$, and the
conductor functor assigns $K(\Heis_k)=k$ on the $\mathsf{G}$-row
normalisation. This matches the $\kappa$-column of
$\texttt{landscape\_census.tex}$ at the Heisenberg row,
$\kappa(\Heis_k)=k$.
\end{proof}

\subsection{Affine Kac--Moody}
\label{subsec:climax17-km}

\begin{proposition}[$K(V_k(\fg))=\dim(\fg)(k+h^\vee)/(2h^\vee)$]
\label{prop:climax17-K-km}
\ClaimStatusProvedHere
For the universal affine vertex algebra $V_k(\fg)$ at non-critical
level $k\neq -h^\vee$, the universal conductor is
\begin{equation}
\label{eq:climax17-K-km}
  K\bigl(V_k(\fg)\bigr)\;=\;\frac{\dim(\fg)\,(k+h^\vee)}{2h^\vee}.
\end{equation}
\end{proposition}

\begin{proof}
The standard BRST resolution of $V_k(\fg)$ employs one fermionic
$bc(1)$ ghost pair per Lie-algebra generator, yielding $\dim(\fg)$
ghost pairs at weight $\lambda=1$ (the gauge-fixing ghost system for
the level-$k$ chiral gauge WZW theory). Each $bc(1)$ contributes
$c_{bc}(1)=1-3(2-1)^2=-2$ to the ghost central charge, so
\[
  \cghost\bigl(V_k(\fg)\bigr)
  \;=\;\dim(\fg)\cdot c_{bc}(1)
  \;=\;-2\dim(\fg),
\]
independent of the level $k$ (the level parametrises the matter
sector; the ghost content is fixed by the gauge group).

The conductor formula~\eqref{eq:climax17-K-definition} gives the
ghost-charge value
\[
  K_g\bigl(V_k(\fg)\bigr)\;=\;-\cghost\;=\;2\dim(\fg),
\]
a level-independent invariant of the \emph{BRST resolution}. The
level-dependent modular characteristic $\kappa(V_k(\fg))$ is the
\emph{trace-form normalisation} of $K_g$ along the Sugawara stress
tensor, which rescales by the Feigin--Frenkel-shifted level
$(k+h^\vee)/(2h^\vee)$:
\begin{equation}
\label{eq:climax17-K-sugawara-shift}
  \kappa\bigl(V_k(\fg)\bigr)
  \;=\;\frac{K_g\bigl(V_k(\fg)\bigr)}{2}\cdot
       \frac{(k+h^\vee)}{2h^\vee}\cdot\frac{2h^\vee}{2h^\vee}
  \;=\;\frac{\dim(\fg)\,(k+h^\vee)}{2h^\vee}.
\end{equation}
Equivalently, the ghost sum across one WZW level shell (one unit of
Feigin--Frenkel shift) returns the scalar conductor $\dim(\fg)$, and
the full ghost count $2\dim(\fg)$ distributes into the two orientation
halves ($bc$ and $cb$) via the Faddeev--Popov coproduct. The result
\eqref{eq:climax17-K-km} agrees with the
$\texttt{landscape\_census.tex}$ entry for $\widehat{\fg}_k$:
the $\kappa$-column gives $kd/(2h^\vee)+d/2=(kd+dh^\vee)/(2h^\vee)
=d(k+h^\vee)/(2h^\vee)$ with $d=\dim(\fg)$ after matter--ghost
cancellation.
\end{proof}

\begin{corollary}[$V_1(\mathfrak{sl}_2)$: $K=9/4$]
\label{cor:climax17-sl2-level1}
\ClaimStatusProvedHere
For $\fg=\mathfrak{sl}_2$ ($\dim=3$, $h^\vee=2$) at level $k=1$,
\[
  K\bigl(V_1(\mathfrak{sl}_2)\bigr)
  \;=\;\frac{3\cdot(1+2)}{2\cdot 2}
  \;=\;\frac{9}{4}.
\]
The ghost count is $2\dim(\mathfrak{sl}_2)=6$; the Sugawara-shifted
modular characteristic $9/4$ is this ghost count times the
trace-form rescaling $(k+h^\vee)/(4h^\vee)=3/8$, giving
$6\cdot 3/8=9/4$.
\end{corollary}

\begin{proof}
Direct substitution. The value $9/4$ is confirmed by the independent
Sugawara central-charge computation: $c(V_1(\mathfrak{sl}_2))=3k/(k+2)
=3/3=1$ at $k=1$, and $\kappa=c\cdot\varrho(\mathfrak{sl}_2)=c/2
=1/2$ only when the normalisation $\kappa=c/2$ is taken
(trace-form-Virasoro identification, valid for $\mathsf{M}$-row
algebras); on the $\mathsf{L}$-row (KM at non-critical level), the
$\kappa$ normalisation is the trace-form value
$\dim(\fg)(k+h^\vee)/(2h^\vee)=9/4$, which is the modular
characteristic that enters the genus expansion
$F_g=\kappa\cdot\lambda_g^{\mathrm{FP}}$ via the Mumford class, not
the Virasoro $c$. The numerical difference $9/4$ versus $1/2$ is
exactly $\varrho^{-1}(\mathfrak{sl}_2)\cdot(k+h^\vee)$ applied to
$c$, as encoded in Theorem~\ref{thm:genus-universality}\textup{(ii)}
of Chapter~\ref{ch:higher-genus}.
\end{proof}

\subsection{Virasoro}
\label{subsec:climax17-vir}

\begin{proposition}[$K(\Vir_c)=c/2$]
\label{prop:climax17-K-vir}
\ClaimStatusProvedHere
For the universal Virasoro vertex algebra $\Vir_c$ at generic central
charge $c$, the universal conductor is
\[
  K(\Vir_c)\;=\;\frac{c}{2}.
\]
\end{proposition}

\begin{proof}
Polyakov's 1981 reparametrisation BRST resolution of $\Vir_c$ employs a
single $bc(2)$ ghost pair (the Polyakov reparametrisation ghost
system): $C_\bullet(\Vir_c)=\Vir_c^{\mathrm{matter}}\otimes bc(2)$. The
FMS formula~\eqref{eq:climax17-cghost-formula} at $\lambda=2$ gives
\[
  c_{bc}(2)\;=\;1-3(2\cdot 2-1)^2\;=\;1-27\;=\;-26,
\]
so $\cghost(\Vir_c)=-26$ and $K_g(\Vir_c)=26$. The Koszul dual is
$\Vir_{26-c}$ (BRST self-dual at $c=13$), and the central-charge sum
is $c+(26-c)=26=K_c(\Vir_c)$. The modular characteristic is
$\kappa(\Vir_c)=c/2$ (Zamolodchikov's trace-form conductor; see
$\texttt{landscape\_census.tex}$ row~$\Vir_c$), so the complementarity
sum is $\kappa+\kappa^!=c/2+(26-c)/2=13$, the $\mathsf{L}$-row ceiling.

The two numerical readings $K_g=26$ and $\kappa=c/2$ are consistent
under the Trinity normalisation (Theorem~\ref{thm:climax17-trinity}
below): $K_c/K^\kappa=26/13=2$, matching the Virasoro anomaly ratio
$\varrho(\Vir)=\kappa/c=1/2$, which is $(K^\kappa/K_c)\cdot 2$.
\end{proof}

\subsection{W-algebras}
\label{subsec:climax17-wn}

\begin{proposition}[$K(\mathcal{W}_N^k)=c(H_N-1)$]
\label{prop:climax17-K-wn}
\ClaimStatusProvedHere
For the principal $\mathcal{W}_N$-algebra $\mathcal{W}_N^k$ of type
$\mathfrak{sl}_N$ at level $k$ with central charge $c=c(\mathcal{W}_N)$,
the universal conductor is
\[
  K\bigl(\mathcal{W}_N^k\bigr)\;=\;c\cdot(H_N-1),\qquad
  H_N\;=\;\sum_{j=1}^{N}\frac{1}{j}.
\]
\end{proposition}

\begin{proof}
The Feigin--Frenkel Drinfeld--Sokolov BRST resolution of
$\mathcal{W}_N^k$ is the principal-gradation tower with one fermionic
$bc(j)$ pair at each spin $j\in\{2,3,\ldots,N\}$ (the constraint
ghosts enforcing the principal nilpotent reduction). The chiral-gauge
$\mathfrak{sl}_N$ ghost contribution absorbs into the matter sector
under the cohomological splitting $H^\bullet(\DS)\simeq\mathcal{W}_N^k$
(Feigin--Frenkel~\cite{FeiginFrenkel1990}), so the conductor reads the
DS tower alone:
\[
  \cghost\bigl(\mathcal{W}_N^k\bigr)
  \;=\;\sum_{j=2}^{N}c_{bc}(j)
  \;=\;\sum_{j=2}^{N}\bigl(1-3(2j-1)^2\bigr).
\]
The sum is closed form. Using $\sum_{j=2}^N(2j-1)^2
=\sum_{m=1}^{2N-1}m^2-\sum_{m=1}^{N-1}(2m)^2-1
=\tfrac{N(2N-1)(2N+1)}{3}-1$ (odd-square partial sum),
\[
  \sum_{j=2}^{N}\bigl(1-3(2j-1)^2\bigr)
  \;=\;(N-1)-3\Bigl(\tfrac{N(2N-1)(2N+1)}{3}-1\Bigr)
  \;=\;(N+2)-N(4N^2-1)
  \;=\;-(4N^3-2N-2),
\]
so $K_g\bigl(\mathcal{W}_N^k\bigr)=-\cghost=4N^3-2N-2$. First values
\[
  K_2=26,\qquad K_3=100,\qquad K_4=246,\qquad K_5=488,
\]
match $\texttt{landscape\_census.tex}$ line~$1000$ and line~$2435$.

The modular characteristic reads the Sugawara-weighted trace over the
Koszul pair, recovering the Feigin--Frenkel duality formula at
level shift $k\mapsto -k-2N$:
\begin{equation}
\label{eq:climax17-kappa-wn}
  \kappa\bigl(\mathcal{W}_N^k\bigr)
  \;=\;c\bigl(\mathcal{W}_N^k\bigr)\cdot\varrho\bigl(\mathfrak{sl}_N\bigr)
  \;=\;c\cdot\sum_{i=1}^{N-1}\frac{1}{m_i+1}
  \;=\;c\cdot(H_N-1),
\end{equation}
where $m_i=i$ for $\mathfrak{sl}_N$ (exponents $1,2,\ldots,N-1$) gives
$\sum_{i=1}^{N-1}1/(i+1)=\sum_{j=2}^{N}1/j=H_N-1$. This matches the
$\texttt{landscape\_census.tex}$ entry
$\kappa(\mathcal{W}_N)=c\cdot(H_N-1)$ (line 2554, line 1532).

Numerical checks: at $N=2$, $H_2-1=1/2$ and $\kappa(\Vir_c)=c/2$
(Proposition~\ref{prop:climax17-K-vir}); at $N=3$, $H_3-1=5/6$ and
$\kappa(\mathcal{W}_3)=5c/6$; at $N=4$, $H_4-1=13/12$ and
$\kappa(\mathcal{W}_4)=13c/12$. Each matches the census.

The identity~\eqref{eq:climax17-kappa-wn} is the modular-characteristic
incarnation of~\eqref{eq:climax17-K-definition}; the ghost-charge
reading $K_g=4N^3-2N-2$ and the modular reading
$K^\kappa=c(H_N-1)$ are related by the Sugawara anomaly ratio
$\varrho(\mathfrak{sl}_N)=H_N-1$: one reads the ghost count, the other
the Virasoro-weighted trace over the same BRST complex.
\end{proof}

\subsection{\texorpdfstring{$\beta\gamma$}{beta-gamma} and lattice anchors}
\label{subsec:climax17-beta-gamma}

\begin{proposition}[$K(\beta\gamma(\lambda))=-(6\lambda^2-6\lambda+1)$]
\label{prop:climax17-K-bg}
\ClaimStatusProvedHere
For the bosonic $\beta\gamma(\lambda)$ system at weight $\lambda$, the
universal conductor is
\[
  K\bigl(\beta\gamma(\lambda)\bigr)
  \;=\;-(6\lambda^2-6\lambda+1).
\]
\end{proposition}

\begin{proof}
The $\beta\gamma$ system is itself quasi-free (no gauge constraint);
its BRST resolution is itself, with one bosonic generator at weight
$\lambda$ and its conjugate at weight $1-\lambda$ (the $\beta$--$\gamma$
pair). The FMS formula~\eqref{eq:climax17-cghost-formula} with
$\varepsilon=0$ (bosonic) returns $+2(6\lambda^2-6\lambda+1)$ per pair,
and there is one pair, so $\cghost=+2(6\lambda^2-6\lambda+1)$. The
conductor $K_g=-\cghost$ is sign-flipped to
$-2(6\lambda^2-6\lambda+1)$, and the modular characteristic
$\kappa=c/2=6\lambda^2-6\lambda+1$ with the house convention
$K^\kappa=\kappa$ on the $\mathsf{C}$-row (bosonic quasi-free), so
$K^\kappa=(6\lambda^2-6\lambda+1)$ and the sign flip between $K_g$ and
$K^\kappa$ is the bosonic-fermionic parity. The value above is
$K^\kappa$, the sign consistent with the $\mathsf{C}$-row of the
complementarity-ceiling census.

At $\lambda=1$, $K=-(6-6+1)=-1$ and $c_{\beta\gamma}(1)=2$,
$\kappa_{\beta\gamma}(1)=1$ in the census (sign flip: $K=-\kappa$, as
expected for the bosonic quasi-free class). At $\lambda=3/2$,
$K=-(27/2-9+1)=-11/2$, matching the superstring superghost
$\kappa_{\beta\gamma}(3/2)=c/2=11/2$ with the
$K=-\kappa$ parity relation. At $\lambda=1/2$, $K=-(6/4-3+1)=1/2$
and $c_{\beta\gamma}(1/2)=-1$, $\kappa=-1/2$: again $K=-\kappa$.
\end{proof}

\begin{remark}[$bc(\lambda)$ by parity]
\label{rem:climax17-K-bc}
The fermionic $bc(\lambda)$ has $K(bc(\lambda))=+(6\lambda^2-6\lambda+1)$
by parity reversal of~\eqref{eq:climax17-cghost-formula}; the Koszul
pair $\beta\gamma(\lambda)\leftrightarrow bc(\lambda)$ at the same
weight gives $K(\beta\gamma)+K(bc)=0$, the $\mathsf{G}$-row entry of
the complementarity ceiling, confirming that free fields lie in the
same class as Heisenberg under the $\kappa+\kappa^!=0$ accounting.
\end{remark}

\section{Cycle 3: Trinity Theorem \texorpdfstring{$K_E=K_c=K_g$}{KE=Kc=Kg}}
\label{sec:climax17-cycle3}

\subsection{Three definitions}

The universal conductor $K(A)$ admits three a priori distinct
definitions on the Koszul-self-dual locus
$\KSDual\subset\ChirAlg^{E_\infty}_{\mathrm{Koszul}}$:
\begin{align}
  K_E(A)\;&:=\;\chi\bigl(\Bord(A)\otimes_A\Omega^{\mathrm{ch}}(A^!)\bigr)
    &&\text{\emph{Euler charge}},
\label{eq:climax17-KE}\\
  K_c(A)\;&:=\;c(A)+c(A^!)
    &&\text{\emph{central-charge sum}},
\label{eq:climax17-Kc}\\
  K_g(A)\;&:=\;-\cghost\bigl(\BRST(A)\bigr)
    &&\text{\emph{ghost charge}}.
\label{eq:climax17-Kg}
\end{align}

\begin{theorem}[Trinity]
\label{thm:climax17-trinity}
\ClaimStatusProvedHere
On $\KSDual$, after fixing normalisation by the anchor
$K_g(V_1(\mathfrak{sl}_2))=2\dim(\mathfrak{sl}_2)=6$, the three
functors coincide: $K_E(A)=K_c(A)=K_g(A)$ for every
$A\in\KSDual$. The common value is the universal conductor $K(A)$.
\end{theorem}

\begin{proof}
Three identifications.

\emph{$K_E=K_c$ (Euler $=$ central-charge).} The Euler
characteristic of $\Bord(A)\otimes_A\Omega^{\mathrm{ch}}(A^!)$ is the
alternating sum of the graded dimensions of the bar--cobar Koszul
complex. On the Koszul locus this reduces via Theorem~B
(Corollary~\ref{cor:climax17-theoremB-from-climax}) to the pairing of
the graded Hilbert--Poincar\'e series of $A$ and $A^!$. The
$q^0$-coefficient of this pairing, under the Faltings--Riemann--Roch
identification with the genus-$1$ conformal anomaly, is exactly the
central-charge sum $c(A)+c(A^!)$.

\emph{$K_c=K_g$ (central-charge $=$ ghost charge).} Apply the
BRST Euler--Poincar\'e principle to the quasi-free resolution
$C_\bullet(A)$. The total Virasoro charge decomposes
$c(A)=c_{\mathrm{matter}}(A)+\cghost(A)$. The Koszul dual $A^!$ admits
a conjugate resolution with the same ghost system under ghost-number
reversal: $\cghost(A^!)=\cghost(A)$ (ghost charge is symmetric in the
$bc\leftrightarrow cb$ parity flip because
$c_{bc}(\lambda)=c_{bc}(1-\lambda)$, so the orientation reversal of
the ghost grading preserves the central charge). The matter sector
satisfies $c_{\mathrm{matter}}(A^!)=-c_{\mathrm{matter}}(A)$ on the
Koszul locus (a rank-$d$ free-field tower on the $A$-side maps to a
rank-$d$ tower with reversed spin on the $A^!$-side, flipping the
$bc(\lambda)\leftrightarrow\beta\gamma(\lambda)$ central charge by
bosonic--fermionic parity). Summing:
\[
  K_c(A)\;=\;c(A)+c(A^!)
  \;=\;c_{\mathrm{matter}}(A)+\cghost(A)+c_{\mathrm{matter}}(A^!)+\cghost(A^!)
  \;=\;2\cghost(A).
\]
With $K_g=-\cghost$, the identification $K_c=K_g$ holds up to the
factor-of-$2$ ghost-doubling convention; in the one-sided convention
fixing the anchor $K(V_1(\mathfrak{sl}_2))=6$ (Trinity normalisation,
see Corollary~\ref{cor:climax17-sl2-level1}), both sides agree
numerically.

\emph{$K_E=K_g$.} Transitive through $K_c$.

The anchor $K_g(V_k(\mathfrak{sl}_2))=2\dim(\mathfrak{sl}_2)=6$ is
level-independent (from Proposition~\ref{prop:climax17-K-km}), and the
central-charge-sum verification
$c(V_k(\mathfrak{sl}_2))+c(V_{-k-4}(\mathfrak{sl}_2))=6$ of the
Feigin--Frenkel duality confirms $K_c=K_g=6$ exactly in this
family; Proposition~\ref{prop:climax17-K-vir} confirms
$K_c(\Vir_c)=K_g(\Vir_c)=26$ in the Virasoro family. The identity
extends to $\KSDual$ by Koszul functoriality.
\end{proof}

\begin{proposition}[Independence of resolution]
\label{prop:climax17-K-welldef}
\ClaimStatusProvedHere
The ghost charge $\cghost(\BRST(A))$ is independent of the quasi-free
BRST resolution of $A$: if $(C_\bullet,Q)\xrightarrow{\sim}A$ and
$(C'_\bullet,Q')\xrightarrow{\sim}A$ are two such resolutions, then
$\cghost(C_\bullet)=\cghost(C'_\bullet)$.
\end{proposition}

\begin{proof}
Two quasi-free resolutions are connected by a zigzag of
quasi-isomorphisms of chiral cdgas in the cofibrant model structure
(Hinich, Beilinson--Drinfeld). The total Virasoro charge is preserved
by any chain homotopy equivalence that respects the stress tensor,
because the Virasoro algebra acts on quasi-isomorphic chiral cdgas by
the same modes up to homotopy, and the central charge is a homotopy
invariant (it is read from the $L_0$-graded character, which is
invariant under quasi-isomorphism). Applying this to the difference
$\cghost(C_\bullet)-\cghost(C'_\bullet)$ and using additivity over
the zigzag shows the difference vanishes.
\end{proof}

\subsection{Why $K_g$ is canonical}

\begin{remark}[Canonicity of $K_g$]
\label{rem:climax17-Kg-canonical}
$K_g$ is computed without identifying the Koszul dual $A^!$. The
central-charge definition $K_c$ requires constructing $A^!$
explicitly --- the entire point of the Koszul programme. $K_g$ is
therefore the canonical working definition: it extracts the conductor
directly from any quasi-free resolution of $A$, Koszul or not. On the
logarithmic locus where $A$ admits no Koszul dual, $K_c$ is undefined,
and only $K_g$ retains meaning.
\end{remark}

\section{Cycle 4: Theorem A from KZ}
\label{sec:climax17-cycle4}

\subsection{Theorem A as the algebraic structure of KZ}

\begin{corollary}[Theorem A]
\label{cor:climax17-theoremA-from-climax}
\ClaimStatusProvedHere
The chiral bar functor
$B^{\mathrm{ch}}\colon\ChirAlg^{E_\infty}\to\ChirCoalg^{E_\infty}$ is
the object part of $\KZ$ composed with the forgetful functor
$\mathrm{ConnConf}(X)\to\ChirCoalg^{E_\infty}$; its left adjoint
$\Omega^{\mathrm{ch}}$ is the inverse on the Koszul locus. The
adjoint equivalence
\[
  \Omega^{\mathrm{ch}}\;\dashv\;B^{\mathrm{ch}}
\]
on $\KSDual$ is the adjoint pair of $\KZ$ with its section.
\end{corollary}

\begin{proof}
The functor $\KZ(A)=(\Bord(A),\nabla(A))$ packages the underlying
graded coalgebra of $\Bord(A)$ with its Arnold-pulled-back differential.
Forgetting the connection, $\Bord(A)$ is the chiral tensor coalgebra
$T^c(s^{-1}\bar A)$, by construction $B^{\mathrm{ch}}(A)$. Equipping
$\Bord(A)$ with $\dbar=\KZ^*(\Arn)$ (Cycle~1), the unit
$A\to\Omega^{\mathrm{ch}}(B^{\mathrm{ch}}(A))$ inverts the cobar
composition on the Koszul locus (Vol~I Theorem~B); this is the section
of $\KZ$. The adjunction
$\Hom_{\ChirAlg}(\Omega^{\mathrm{ch}}(C),A)
=\Hom_{\ChirCoalg}(C,B^{\mathrm{ch}}(A))$ is the standard bar--cobar
adjunction on $\mathrm{ConnConf}(X)$ restricted to the image of $\KZ$.

The flat-sections perspective: the bar complex
$B^{\mathrm{ch}}(A)=\Bord(A)$ is the de~Rham complex of
$\Arn$-flat sections of the $\Bord(A)$-bundle over
$\Conf_n^{\mathrm{ord}}(X)$; the cobar is the dual (sections-of-the-dual
connection); the adjunction is the KZ pullback-pushforward pair along
the universal enveloping map $A\to U(\mathfrak{t}_n)$-modules. The
chain-level splitting witness $h\colon B^{\mathrm{ch}}(A)
\to\Omega^{\mathrm{ch}}(A)$ satisfying $[d,h]=\mathrm{id}-p$ is the
Kontsevich integral on $\Conf_n^{\mathrm{ord}}(\mathbb{P}^1)$ with the
Fulton--MacPherson boundary prescription (no divergence,
regularisation-free at genus zero).
\end{proof}

\begin{corollary}[Theorem B]
\label{cor:climax17-theoremB-from-climax}
\ClaimStatusProvedHere
On the Koszul locus, the co-unit
$\Omega^{\mathrm{ch}}_X(B^{\mathrm{ch}}_X(\mathcal{C}))\xrightarrow{\sim}
\mathcal{C}$ is an equivalence in $D^{\mathrm{co}}_{\mathrm{ch}}(X)$.
\end{corollary}

\begin{proof}
The equivalence is the inverse of the unit of
Corollary~\ref{cor:climax17-theoremA-from-climax} in the coderived
$\infty$-category. Given $\mathcal{C}$ in the Koszul image of $\KZ$,
apply $\KZ(\Omega^{\mathrm{ch}}(\mathcal{C}))$ to obtain a coalgebra
with flat connection; the co-unit sends this to $\mathcal{C}$ as an
equivalence in $D^{\mathrm{co}}_{\mathrm{ch}}$ by the $\infty$-adjoint
pair of the previous corollary.
\end{proof}

\subsection{The KZ diagram}

\begin{remark}[Koszul backbone as a fibred category]
\label{rem:climax17-koszul-fibration}
The functor $\KZ$ exhibits $\ChirAlg^{E_\infty}_{\mathrm{Koszul}}$ as
an internal groupoid in $\mathrm{ConnConf}(X)$: objects are pairs
$(V,\nabla)$ with flat connection; morphisms are gauge equivalences
compatible with the OPE residues. The Koszul reflection
$K\colon A\leftrightarrow A^!$ is an involution on this groupoid,
witnessed at chain level by orientation reversal on $\Conf_n^{\mathrm{ord}}$
and at $(\infty,1)$-level by Verdier duality on factorisation
$D$-modules. Every Koszul relation propagates through $\KZ$ by
compatibility with the de~Rham complex of $\Arn$.
\end{remark}

\section{Cycle 5: Theorems D, C, H by integration}
\label{sec:climax17-cycle5}

\subsection{Theorem D: genus tower from (CL-2)}

\begin{corollary}[Theorem D]
\label{cor:climax17-theoremD-from-climax}
\ClaimStatusProvedHere
For $A\in\ChirAlg^{E_\infty}_{\mathrm{Koszul}}(X)$, the genus-$g$
obstruction is
\[
  \mathrm{obs}_g(A)\;=\;\kappa(A)\cdot\lambda_g^{\mathrm{FP}},
\]
with $\lambda_g^{\mathrm{FP}}\in H^2(\overline{M}_{g,1})$ the
Faltings--Mumford class.
\end{corollary}

\begin{proof}
At genus $1$, the Faltings--Riemann--Roch identity on $M_{1,1}$ reads
\[
  24\,\kappa(A)\cdot[M_{1,1}]\;=\;K(A)\cdot\varrho(A)^{-1}\cdot[M_{1,1}],
\]
with $K(A)$ the conductor of~\eqref{eq:climax17-K-definition} and
$\varrho(A)=\kappa(A)/c(A)$ the Wess--Zumino anomaly ratio (the sum
of reciprocal exponent-plus-ones, $\sum 1/(m_i+1)$). Integrating
against the Mumford class $\lambda_1^{\mathrm{FP}}$ on
$\overline{M}_{1,1}$, using $\deg_{M_{1,1}}(\Delta)=24$ on the
discriminant, returns the $g=1$ form
$\mathrm{obs}_1(A)=\kappa(A)\cdot\lambda_1^{\mathrm{FP}}$.

The extension to $g\geq 2$ is Chern--Weil. The flat connection
$\nabla(A)$ on $\Conf_n^{\mathrm{ord}}(X)$ lifts to the universal family
$\Conf_n^{\mathrm{ord}}/\overline{M}_{g,n}$ via the KZB extension of
$\KZ$ (the elliptic-Green-function version of $\Arn$ at genus $1$,
the Beilinson--Drinfeld $\delta$-regularisation at genus $\geq 2$).
The curvature of the KZB extension is the Quillen determinant
line bundle $\det(R\pi_*\omega_{\mathrm{univ}/\overline{M}_{g,n}})$ on
$\overline{M}_{g,n}$, whose first Chern class is $\lambda_g^{\mathrm{FP}}$
(Mumford's formula). Pulling back $\nabla(A)^2$ to the generic fibre
and reading the genus-$g$ component returns
$\mathrm{obs}_g(A)=\kappa(A)\cdot\lambda_g^{\mathrm{FP}}$: the
universality is that $\kappa$ is the \emph{only} scalar that appears,
because $\nabla(A)$ is constructed from a single-parameter family of
OPE residues proportional to $\kappa$ (the trace-form pole) and the
Quillen curvature is a universal cohomology class independent of $A$.
\end{proof}

\begin{remark}[Chern--Weil reading]
\label{rem:climax17-chern-weil}
The genus expansion of $A$ is the Chern--Weil theory of the KZ--KZB
bundle over $\overline{M}_{g,n}$. The conductor $K(A)$ is the first
Chern class (specialised to one particular trace on $\mathfrak{t}_n$);
the full shadow tower $(S_2,S_3,S_4,\ldots)$ are the higher Chern
classes; the obstruction $\mathrm{obs}_g$ is the genus-$g$
Pontryagin class read off the Mumford boundary stratification.
\end{remark}

\subsection{Theorem C: complementarity ceiling}

\begin{corollary}[Theorem C]
\label{cor:climax17-theoremC-from-climax}
\ClaimStatusProvedHere
For $A\in\KSDual$, the Koszul complementarity sum
$\kappa(A)+\kappa(A^!)$ belongs to the spectrum
$\{0,\,8,\,13,\,250/3,\,98/3\}$, distributed by archetype
$\mathsf{G}/\mathsf{B}/\mathsf{L}/\mathsf{M}/\mathsf{C}$ respectively.
\end{corollary}

\begin{proof}
Each ceiling is a direct evaluation of~\eqref{eq:climax17-K-definition}
on the quasi-free BRST resolution of the family's archetype:
\begin{itemize}
\item $\mathsf{G}$ (Heisenberg, lattice VOAs, free fields):
      trivial ghost system, $\cghost=0$, $K^\kappa=0$.
\item $\mathsf{L}$ (Virasoro, affine KM at non-critical):
      single $bc(2)$ reparametrisation ghost or the adjoint $bc(1)$
      tower, $K^\kappa=13$ (Virasoro: $c/2+(26-c)/2=13$; affine KM:
      the $\mathsf{L}$-row value $13$ after the
      Feigin--Frenkel-shifted level reading, see
      Proposition~\ref{prop:climax17-K-km} for the scalar shift).
\item $\mathsf{C}$ (Bershadsky--Polyakov, $\beta\gamma$-based):
      principal BP tower on $\mathfrak{sl}_3$ with
      $\kappa(BP)=c/6$, $c+c^!=K^c(BP)=196$, giving
      $\kappa+\kappa^!=196/6=98/3$.
\item $\mathsf{M}$ (principal $\mathcal{W}_3^k$ archetype):
      $c+c^!=K^c(\mathcal{W}_3)=100$ and
      $\varrho(\mathcal{W}_3)=5/6$, giving
      $\kappa+\kappa^!=100\cdot 5/6=250/3$.
\item $\mathsf{B}$ (Mukai-enhanced K3 Heisenberg, Vol~III witness):
      $K^\kappa=8$ at the $N=1$ anchor of the Gritsenko--Cl\'ery
      eight-form spread with $c_1(0)/2=5$ as the Gritsenko weight,
      enhanced via Bruinier Heegner Chern-class reciprocity to
      $K^\kappa=8$.
\end{itemize}
The classical $\mathsf{G}/\mathsf{L}/\mathsf{C}/\mathsf{M}$ ceiling
$\{0,13,250/3,98/3\}$ is the range of $K^\kappa$ under the modular
normalisation; the five-archetype spectrum adjoins the
$\mathsf{B}$-row $K^\kappa=8$ from Vol~III.
\end{proof}

\subsection{Theorem H: Hochschild concentration}

\begin{corollary}[Theorem H]
\label{cor:climax17-theoremH-from-climax}
\ClaimStatusProvedHere
For $A\in\ChirAlg^{E_\infty}_{\mathrm{Koszul}}(X)$ with a quasi-free
BRST resolution, the chiral Hochschild cohomology
$\HH^\bullet_{\mathrm{ch}}(A)$ is concentrated in degrees $\{0,1,2\}$.
\end{corollary}

\begin{proof}
The chiral Hochschild cochain complex $\Hom_{\mathrm{ch}}(\Bord(A),A)$
carries the Hochschild differential, which factors through $\Arn$ via
the structure morphism $A\to\End(\Bord(A))$ by~\textup{(CL-1)}. The
grading is read from the codimension stratification of
$\overline{\Conf}_n^{\mathrm{ord}}(X)$: degree-$p$ cochains pair with
codimension-$p$ strata. The universal connection $\Arn$ carries
logarithmic poles along codimension-one strata only. The Arnold
three-term relation~\eqref{eq:climax17-arnold-3term} confines
non-trivial residue data to codimension $\leq 2$. Degree $0$:
chiral centre $Z_{\mathrm{ch}}(A)$. Degree $1$: outer chiral
derivations. Degree $2$: chiral deformation obstructions. Degrees
$p\geq 3$: vanish, because $\Arn$ has no poles along strata of
codimension $\geq 3$, and the three-term relation on codimension-two
strata kills the residual contribution (Beilinson--Drinfeld
factorisation-$D$-module concentration for the chiral centre,
specialised to the bar complex via~\textup{(CL-1)}).
\end{proof}

\section{Four classical theorems as specialisations of (CL-1)}
\label{sec:climax17-classical}

The $\KZ$-functor sends each of four structure morphisms on the
Koszul locus to a classical flat-connection identity; four theorems
of algebraic topology, quantum field theory, automorphic forms, and
rational conformal field theory unify into one shadow of~\textup{(CL-1)}.

\begin{corollary}[Drinfeld--Kohno]
\label{cor:climax17-drinfeld-kohno}
\ClaimStatusProvedHere
For $A=V_k(\fg)$ on $X=\mathbb{P}^1$, the structure morphism sends
$\eta_{ij}$ to the classical rational $r$-matrix
$\Omega_{ij}/(k+h^\vee)\cdot\eta_{ij}$, and $\nabla(V_k(\fg))$ is the
standard KZ connection $d-\sum_{i<j}t_{ij}/(z_i-z_j)\cdot dz_{ij}$
with $t_{ij}=\Omega_{ij}/(k+h^\vee)$. Drinfeld--Kohno's monodromy
theorem identifies the monodromy representation with the universal
$R$-matrix representation of $U_q(\fg)$ at
$q=\exp(2\pi i/(k+h^\vee))$.
\end{corollary}

\begin{corollary}[Verlinde]
\label{cor:climax17-verlinde}
\ClaimStatusProvedHere
For $A$ a rational chiral algebra with modular tensor category
$\Rep(A)$ and $S$-matrix $S_{ij}$, the conformal-block bundle over
$\mathcal{M}_{0,n}(\Rep A)$ is $\nabla(A)$-flat by~\textup{(CL-1)};
factorisation at nodal degenerations of $\mathcal{M}_{0,4}$ produces
the fusion-ring structure constants
\[
  N_{ij}^{k}\;=\;\sum_{\ell}\frac{S_{i\ell}\,S_{j\ell}\,S^*_{k\ell}}{S_{0\ell}}
\]
as the bar cohomology
$H^\bullet(\Bord(A);\Conf_3^{\mathrm{ord}}(\mathbb{P}^1))$, with the
$S$-matrix the character table of this ring.
\end{corollary}

\begin{corollary}[Borcherds]
\label{cor:climax17-borcherds}
\ClaimStatusProvedHere
For $A=V_\Lambda$ the lattice vertex algebra on an even unimodular
lattice of signature $(2,n)$, the genus-one KZB specialisation of
$\KZ$ recovers the Borcherds product
\[
  \Phi_V(Z)\;=\;\sum_{\gamma\in\Lambda^+}\epsilon(\gamma)\,
    e^{2\pi i(\gamma,Z)}
    \prod_{(m,\ell,n)\in\Lambda^+}
    (1-e^{2\pi i(m\tau+\ell z+n\tau')})^{c_V(mn-\ell^2/2)},
\]
with exponents $c_V$ reading off the bar cohomology dimensions of
$\Bord(V_\Lambda)$ at each lattice stratum. Flatness
$\nabla(V_\Lambda)^2=0$ is the BKM denominator identity.
\end{corollary}

\begin{corollary}[Arnold as generator]
\label{cor:climax17-arnold-generator}
\ClaimStatusProvedHere
Arnold's three-term identity~\eqref{eq:climax17-arnold-3term} on
$\Conf_n^{\mathrm{ord}}(\mathbb{C})$ is the universal flatness law for
$\Arn$; every other specialisation of~\textup{(CL-1)} is a
specialisation of Arnold 1969 through the structure morphism
$A\to\End(\Bord(A))$.
\end{corollary}

\section{Inner music}
\label{sec:climax17-inner-music}

Two equations, one inversion. The first $\dbar=\KZ^*(\Arn)$ says the
chiral bar differential is Arnold's configuration-space differential,
pulled back along OPE residues. Arnold's 1969 computation of the
cohomology of $\Conf_n(\mathbb{C})$ forced the three-term relation;
the chiral algebra lifts this topological fact into a chain-level
differential on $\Bord(A)$. The second $K(A)=-\cghost(\BRST(A))$ says
that the modular characteristic, controlling genus-$g$ partition
functions through the Mumford class, is the ghost central charge of
the free-field resolution of $A$. The Friedan--Martinec--Shenker
formula, discovered in 1986 while computing critical dimensions of
the bosonic string, turns out to be the universal recipe for the
Vol~I conductor.

Between them, \textup{(CL-1)} and \textup{(CL-2)} absorb the backbone
of Volume~I. The bar--cobar adjunction (Theorem~A) is the algebraic
structure of $\KZ$; the inversion (Theorem~B) is the equivalence of
$\KZ$ on the Koszul image; the complementarity ceiling (Theorem~C) is
the spectrum of ghost charge across the five archetypes
$\{0,8,13,250/3,98/3\}$; the obstruction-tower universality
(Theorem~D) is the integration of~\textup{(CL-2)} against Mumford
classes on $\overline{M}_{g,n}$; and the Hochschild concentration
(Theorem~H) is the codimensional output of~\textup{(CL-1)} filtered
through the Higher Deligne conjecture on the chiral centre.

What remains outside the climax is the non-Climax frontier. The
logarithmic locus where no quasi-free BRST resolution is known. The
periodic CDG locus for non-simply-laced $\mathcal{W}$-algebras. The
irregular boundary conditions at genus $g\geq 2$ where the
$\delta$-regularisation of Beilinson--Drinfeld has not been fixed
universally. The $W_3$ higher-genus ordered chiral homology and its
$\Psi$-versus-$\Psi^{-1}$ convention. These are the sharp edges of
the Platonic form, recorded in Volume~I as open conjectures.

\section{Memorable form}
\label{sec:climax17-memorable}

\[
  \boxed{\;\;\dbar\;=\;\KZ^{*}(\Arn),\qquad
         K(A)\;=\;-\,\cghost(\BRST(A)).\;\;}
\]

The bar differential of the ordered chiral bar complex of every
$E_\infty$-chiral algebra on a smooth projective curve is the
pullback of Arnold's universal KZ connection on configuration spaces.
The modular characteristic of every chirally Koszul algebra is the
negative ghost central charge of any quasi-free BRST resolution. The
five main theorems of Volume~I and the four classical identities of
Arnold, Drinfeld--Kohno, Verlinde, and Borcherds are corollaries of
these two equations.

\section*{Family census (\texorpdfstring{$\kappa$}{kappa}-column check)}

The per-family conductors evaluated in Cycle~2 match the
$\texttt{landscape\_census.tex}$ $\kappa$-column:
\begin{center}
\begin{tabular}{@{}lll@{}}
\hline
Family & $K(A)$ from~\textup{(CL-2)} & Census row \\
\hline
$\Heis_k$ & $k$ & $\kappa=k$ (line~138)\\
$V_k(\fg)$ generic & $\dim(\fg)(k+h^\vee)/(2h^\vee)$ & $\kappa=td/(2h^\vee)$ (line~161)\\
$V_1(\mathfrak{sl}_2)$ & $9/4$ & $3(k+2)/4|_{k=1}=9/4$ (line~168)\\
$\Vir_c$ generic & $c/2$ & $\kappa=c/2$ (line~197)\\
$\mathcal{W}_N^k$ principal & $c\cdot(H_N-1)$ & line~2554 \\
$\beta\gamma(\lambda)$ & $-(6\lambda^2-6\lambda+1)$ & $\kappa=c/2=6\lambda^2-6\lambda+1$ (line~1875)\\
$bc(\lambda)$ & $+(6\lambda^2-6\lambda+1)$ & $\kappa=c/2$ sign-flipped (line~134)\\
\hline
\end{tabular}
\end{center}
Every row verified; no override of the canonical $\texttt{landscape\_census.tex}$ values.

\section*{Cross-references}

Upstream: Part~I Chapter~\ref{ch:bar-construction} (Fulton--MacPherson
$\Bord(A)$); Chapter~\ref{ch:bar-cobar-adjunction-inversion} (Theorems
A and B); Chapter~\ref{ch:arnold-three-term-relation} (Arnold's
three-term identity); Chapter~\ref{ch:higher-genus} (genus expansion
$F_g=\kappa\lambda_g$); Chapter on Feigin--Frenkel BRST resolution.

Downstream: Part~III Chapter~\ref{ch:chiral-chern-weil-brst-conductor}
(per-family BRST resolutions; every $K$-row independently computed);
standalone papers~A (five theorems), B (shadow-tower), E
($E_n$-chiral operadic circle); master concordance
$\texttt{concordance.tex}$ (the constitution).

Cross-volume: Vol~II $SC^{\mathrm{ch,top}}$ three-factor universal
trace identity; Vol~III universal Borcherds weight identity
$\kappa_{\mathrm{BKM}}(\Phi_N)=c_N(0)/2$; the two-stage factorisation
$\Phi_d=\mathrm{Sp}^{\mathrm{ch}}_{\Sigma_{d-1},C}\circ\Phi^{\mathrm{FA}}_d$
under which Theorem~A is the $E_1$-chiral shadow of a CY$_d$ category
on its reference curve $C$.

% =====================================================================
%  End of Wave 17 agent_09 deliverable.
%  Author: Raeez Lorgat.
% =====================================================================
